% Pregunta 405 (2011-octubre-OEP2011, n.º 26): teorema de Rybczynski en la caja de los factores; al aumentar el capital con precios constantes, sube la producción del bien intensivo en capital (X) y baja la del intensivo en trabajo (Y). Autor: Víctor Gutiérrez Marcos
% A escala: L̄ = 6, K̄ = 4; K_X/L_X = 1,5 (X intensivo en capital), K_Y/L_Y = 0,3.
%   A = (1,8333; 2,75): L_Y = 4,1667, K_Y = 1,25.
% K̄ sube a 5,4 (+35 %): 1,5 L_X + 0,3 (6 - L_X) = 5,4 da A' = (3; 4,5):
%   X sube un 63,6 % (más que K̄) e Y baja un 28 % (L_Y de 4,1667 a 3).
\documentclass[tikz,border=6pt]{standalone}
\input{tikz-tcee}
\begin{document}
\begin{tikzpicture}[tcee,x=1cm,y=1cm,
    llave/.style={decorate,decoration={brace,amplitude=5pt},thin,tceegris}]
  % Caja inicial L̄ × K̄
  \draw[thick] (0,0) rectangle (6,4);
  \node[below left,text=tceemorado] at (0,0) {$O_X$};
  \node[right,text=tceeverde] at (6,4) {$O_Y$};
  \draw[llave] (6,-0.3) -- (0,-0.3) node[midway,below=5pt,text=black] {$\bar L$ (trabajo)};
  \draw[llave] (-0.3,0) -- (-0.3,4) node[midway,left=5pt,text=black] {$\bar K$};
  % Aumenta K̄: la caja se alarga hacia arriba y O_Y sube a O_Y'
  \draw[curvanueva] (0,4) -- (0,5.4) -- (6,5.4) -- (6,4);
  \node[above right,text=tceerojo] at (6,5.4) {$O_Y'$};
  \draw[llave,tceerojo] (-1.25,0) -- (-1.25,5.4) node[midway,left=5pt,text=tceerojo] {$\bar K'$};
  \draw[flecha] (6.35,4.1) -- (6.35,5.3);
  % Rayo de X desde O_X (empinado: X intensivo en capital)
  \draw[curva1] (0,0) -- (3.3,4.95);
  \node[etiqueta,text=tceemorado,anchor=west] at (0.6,0.3) {rayo de $X$};
  % Rayo de Y desde O_Y (poco inclinado) y punto A
  \draw[curva3] (6,4) -- (1.2,2.56);
  \node[etiqueta,text=tceeverde,anchor=north] at (4.4,3.4) {rayo de $Y$};
  \node[punto,label={[etiqueta]above left:$A$}] (A) at (1.8333,2.75) {};
  % Rayo de Y desde O_Y' (paralelo: precios de los factores constantes) y A'
  \draw[curvanueva,desplazada] (6,5.4) -- (2.4,4.32);
  \node[punto,label={[etiqueta]above left:$A'$}] (A2) at (3,4.5) {};
  % Rótulo
  \node[etiqueta,align=center,text=tceerojo] at (3.55,1.25)
    {$X$ aumenta: $A'$ más lejos de $O_X$\\$Y$ disminuye: $A'$ más cerca de $O_Y'$};
\end{tikzpicture}
\end{document}
